\chapter{Divisors and Abel maps}
\label{pa-ch-divisors}

\begin{definition}[Symmetric powers]
  \label{pa-def-symmetric-power}
  \dcref{III.6,custom}
  \cite{milne-av,kleiman-picard}
  For \(n\ge0\), \(C^{(n)}\) is the quotient of \(C^n\) by the symmetric group.
  Its geometric points are effective divisors of degree \(n\).
\end{definition}

The Hilbert scheme of length-\(n\) subschemes of a smooth curve represents
this functor.  Its local coordinates are the coefficients of a monic
polynomial, so it is smooth of relative dimension \(n\).

\begin{proposition}[Effective divisor map]
  \label{pa-prop-effective-map}
  \dcref{3.7}
  \uses{pa-def-symmetric-power,pa-def-picard}
  \cite[Thm.~3.7]{kleiman-picard}
  The assignment \(D\mapsto\mathcal O_C(D)\) gives a morphism
  \[
    C^{(n)}\longrightarrow\operatorname{Pic}^n_{C/k}
  \]
  whose geometric fibres are complete linear systems.
\end{proposition}

\begin{proof}
  A finite flat effective divisor is Cartier on the smooth relative curve,
  hence determines an invertible sheaf of degree \(n\).  The construction is
  compatible with pullback and factors through the etale Picard sheaf.  Two
  divisors have the same image exactly when their difference is principal,
  which identifies the geometric fibre with the complete linear system.
\end{proof}

\begin{definition}[Abel--Jacobi map]
  \label{pa-def-abel-jacobi}
  \dcref{III.6,custom}
  \uses{pa-def-picard}
  \cite{milne-av,kleiman-picard}
  For \(P_0\in C(k)\), set
  \[
    \iota_{P_0}(Q)=[\mathcal O_C(Q-P_0)].
  \]
  The rigidified diagonal line bundle makes
  \(\iota_{P_0}:C\to\operatorname{Pic}^0_{C/k}\) a morphism with
  \(\iota_{P_0}(P_0)=0\).
\end{definition}

\begin{proof}
  The diagonal divisor on \(C\times C\), corrected by the two fibres through
  \(P_0\), is rigidified along \(P_0\).  Its restriction over \(Q\) is
  \(\mathcal O_C(Q-P_0)\).  The universal property of the relative Picard
  sheaf therefore gives the stated morphism and sends \(P_0\) to the identity.
\end{proof}
